\begin{textsample}{NEG Dim 4 – pb – Score -1.00 – pb/pb\_ef\_4.txt}  \label{ex:f4_neg_130}
A área de Linguagens reúne os componentes curriculares : Língua Portuguesa, Arte, Educação Física e, no Ensino Fundamental – Anos Finais, Língua Inglesa.
Assim, cada componente curricular da área evidencia linguagens que passam a ter status próprio de objetos de conhecimento escolar na presente proposta curricular, com a intenção de possibilitar aos alunos “ participar de práticas de linguagem diversificadas, que lhes permitam ampliar suas capacidades expressivas em manifestações artísticas, corporais e linguísticas, como também seus conhecimentos sobre essas linguagens ” ( BRASIL, 2017, p. 61 ).
A seguir são apresentadas as propostas específicas para cada um dos componentes da área. 5.1 - Língua Portuguesa 5.1.1 Introdução O Currículo de Língua Portuguesa do Ensino Fundamental visa garantir os direitos de aprendizagem e desenvolvimento do aluno, \textbf{consubstanciados} nas competências gerais e específicas bem como nas habilidades delineadas na BNCC, considerando as práticas de linguagem situadas nos campos artístico/literário, de estudo e pesquisa, jornalístico/midiático, de atuação da vida pública e pessoal, cidadã e investigativa, na interface de 04 ( quatro ) eixos : Oralidade, Leitura/escuta/interpretação, Produção ( escrita e multissemiótica ) e Análise linguística/semiótica ( que envolve conhecimentos linguísticos – sobre o sistema de escrita, o sistema da língua e a norma padrão – textuais, discursivos e sobre os modos de organização e os elementos de outras semioses ).
A oralidade, contrariamente à escrita, que historicamente tem status garantido na escola, vem se fixando como eixo norteador do ensino de língua/linguagem, principalmente depois das orientações dos Parâmetros Curriculares Nacionais ( PCN ).
Possui representações informais e formais, com normas e valores sociais inerentes a cada uma.
Assim sendo, a escola deve necessariamente colocar em seu plano de trabalho ações didáticas que orientem o sujeito aprendente a empregar a linguagem oral adequadamente nas mais diversas situações comunicativas.
Como toda comunicação verbalizada concretiza -se em textos, é importante para a prática pedagógica a seleção de gêneros textuais orais diversos, associados à situação comunicativa específica, para assegurar não só o reconhecimento desses gêneros, mas, paulatinamente, fazer com que o aluno se aproprie de práticas orais que contemplem desde situações do cotidiano até as que exigem mais formalidade.
Para tanto, torna -se crucial possibilitar leituras e vivências com a modalidade oral, desde a Educação Infantil, que possibilitem a locução e interlocução dos alunos, o respeito ao turno da fala, o planejamento da fala, além de aspectos paralinguísticos próprios dessa modalidade.
Esses saberes que deverão ser progressivamente ampliados no ensino fundamental, por meio da leitura em voz alta ( exercício para pronúncia clara dos fonemas ), apresentação de argumentos, vivências com gêneros jornalísticos televisivos e/ou radiofônicos ou com as novas mídias de comunicação, além de apresentação de seminários, entre outros, são maneiras relevantes de preparar os alunos para práticas escolares e para a vida em sociedade, que lhes permitam diferenciar as comunicações nos contextos sociais particulares ( família, amigos, trabalho ) dos usos com mais formalidades, como conferências, discursos políticos, religiosos, científicos etc.
São formas de preparar o aluno para o pleno exercício da cidadania, conscientizando -o sobre a diversidade linguística/textual/discursiva, coerente com uma prática situada socioculturalmente, o que implica ainda desenvolver a maturidade para a convivência social respeitosa, que valorize o outro, a esfera social e a si mesmo.
A leitura, por sua vez, permite ampliar as possibilidades tanto de se aprender nas demais áreas de estudo escolar quanto de se adquirir autonomia para construção dos conhecimentos fora da escola.
Ela é, inquestionavelmente, o caminho para a formação do cidadão crítico, reflexivo e, sobretudo, autônomo no seu agir.
Isso porque amplia os horizontes do conhecimento, favorece o acesso a novas ideias, permite o diálogo, promove a mudança de percepções e de concepções, contribui para a intervenção sobre o mundo, incluindo o homem no exercício da cidadania e na tomada de decisões da própria vida.
Daí a sua importância em ocupar lugar de destaque na escola, pois é a prática sistematizada de estudo de gêneros e temas diversos que possibilita, ao aluno, perceber o seu entorno e se perceber como participante na interação com o mundo, exercendo papéis que efetivam a construção de uma sociedade melhor.
Isso é possível porque o trabalho com gêneros textuais produzidos em situações comunicativas diferentes e com o emprego de linguagens diversas conduz o leitor a refletir sobre as condições de produção : quem escreveu, por que, quando, com que finalidade, para que público, onde circulou.
São conhecimentos que, quando se associam aos conhecimentos prévios do leitor, contribuem para a construção de sentido.
Mas, para isso, é preciso que o professor conduza o aluno por esse caminho, fazendo -o vencer o nível da decodificação e atingir o objetivo maior da interpretação do texto e da percepção das intenções com que foi produzido.
Tal ideia perpassa os textos do universo escolar, do universo literário e do mundo prático, contribuindo para a formação integral do aluno, mas sem que haja a supervalorização de um em detrimento de outro. > Ler é uma forma de saber o que se passa, o que se pensa, o que se diz, é uma forma de ficar inteirado acerca do que vai pelo mundo, acerca do que vai povoando a cabeça e coração dos pensadores, dos formadores de opinião, dos cientistas, dos poetas.
É uma forma de saber acerca das descobertas que foram feitas ou das hipóteses que estão sendo testadas, ou dos planos e projetos em andamento ( ANTUNES, 2009, p. 195 ). > [... ] a oralidade e escrita não estão em competição.
Cada uma tem a sua história e seu papel na sociedade ( MARCUSCHI, L.
Antônio, 2005, p. 15 ).
Os PCN apresentam orientações específicas para o ensino da produção textual : > “ o trabalho com produção de textos tem como finalidade formar produtores competentes e capazes de produzir textos coerentes, coesos e eficazes ” ( BRASIL, 1997, p. 47 ). > A leitura não depende apenas do contexto linguístico do texto, mas também do contexto extralinguístico de sua produção e circulação ( ANTUNES, 2003, p. 77 ).
Na esfera escolar, a leitura de textos didáticos, resenhas, artigos de divulgação científica, por exemplo, são importantes porque favorecem o aprendizado escolar, exigem a capacidade de associar temas interdisciplinares.
A leitura de textos literários, por sua vez, apresenta ao aluno uma forma artística e estética de lidar com questões reais, dialogando, portanto, com a vida humana.
Já a leitura de gêneros do mundo prático, tais como notícias e propagandas, permite uma visão da construção da realidade.
As três esferas, portanto, são complementares e dão sentido ao mundo, indicando caminhos para a participação em diferentes campos da atividade humana.
A produção textual deve ser orientada pela concepção dialógica da linguagem, cujo escopo é a unicidade do ato comunicativo, que se organiza por gêneros textuais/discursivos, que são, por sua vez, determinados pelas condições de produção que envolvem os interlocutores e a situação comunicativa.
Em vista disso, a noção de gênero textual é fundamental para a eficácia do ensino e da aprendizagem da leitura e produção textual, a fim de assegurar um aprendizado significativo e apropriado para o desenvolvimento de sujeitos leitores/escritores.
Por isso mesmo, as atividades de produção textual devem favorecer o acesso dos alunos aos mais variados gêneros empregados nas diversas práticas sociais, para que eles possam organizar seus discursos conscientes das suas intenções, de seus motivos e de suas razões ( BRONCKART, 2006 ).
Devem ainda servir para ampliar as práticas sociais e as de letramento, de forma crítica e reflexiva, contribuindo, assim, para o exercício da cidadania plena e autônoma, nos espaços particulares e/ou coletivos.
Para isso, é importante que a prática pedagógica valorize a refacção/reescrita de um texto, observando a adequação dos aspectos linguísticos, discursivos e composicionais de um gênero no interior de um texto ; ou a retextualização, observando a transformação de um gênero em outro, seja de um texto oral para outro oral ; seja de um oral para um escrito, ou vice-versa ; seja de um escrito para outro escrito ; ou até mesmo empregando configurações textuais, tais como : do não verbal para o verbal ou vice-versa ; do texto multimodal para o oral, ou para o escrito, ou para o gestual etc. > [... ] a reescrita é atividade na qual, através do refinamento dos parâmetros discursivos, textuais e linguísticos que norteiam a produção original, materializa -se uma nova versão do texto.
Já na retextualização, tal como entendida aqui, opera -se, fundamentalmente, com novos parâmetros de ação da linguagem, porque se produz novo texto : trata -se, além de redimensionar as projeções de imagem dos interlocutores, de seus papéis sociais e comunicativos, dos conhecimentos partilhados, assim como de motivações e intenções, de espaço e tempo de produção/recepção, de atribuir novo propósito à produção linguageira ( MATÊNCIO, 2002, p. 112 ) ( Grifo nosso ).
Finalmente, a análise linguística/semiótica tem no texto seu ponto de partida e de chegada, e todo o percurso deve salientar a triangulação uso-reflexão-uso, conduzindo à percepção de como a língua/linguagem funcionam para a construção de sentidos, ou seja, permitindo explicar o que o texto diz e de que recursos da língua ele faz uso para dizer o que diz.
Logo, é uma análise feita a partir de elementos que são importantes para a compreensão do texto, por isso, ela não pode ser isolada, mas contextualizada e reflexiva, o que implica estar sempre articulada com as práticas de leitura, de oralidade e de produção textual, até porque as aulas de Língua Portuguesa devem servir para pensar que tanto as escolhas lexicais quanto a organização linguística/discursiva influenciam para a construção de sentido nas práticas interativas no meio social.
Sendo assim, o estudo da gramática deve permitir ao aluno reconhecer e validar as várias formas de dizer e a sua adequação às diversas realidades sociocomunicativas. > O enfoque deve ser essencialmente semântico-pragmático-discursivo : as reflexões sobre os aspectos especificamente gramaticais precisam ser lançadas contra esse pano de fundo semântico-pragmático-discursivo, de modo a conscientizar o aprendiz de que os recursos disponíveis na língua são ativados essencialmente para a produção de sentido e a interação social. ( BAGNO, 2011, p. 20 ).
O professor deve levar para sala de aula uma gramática : - que seja relevante, selecionando noções e regras gramaticais que sejam úteis e aplicáveis à compreensão e aos usos sociais da língua ; - que seja funcional, propondo uma gramática referendada no funcionamento efetivo da língua ; - contextualizada, que esteja incluída na interação verbal ; - que traga algum tipo de interesse, sendo estimulante e desafiador ; - que liberte, incentivando a fluência linguística e neutralizando a postura prescritiva e corretiva nas produções dos alunos ; - que prevê mais de uma norma, caracterizando adequadamente a norma padrão como a privilegiada, mas não como a única certa ; - que é da língua e das pessoas, discernindo o que é significativo para a experiência humana da interação verbal. ( ANTUNES, 2009 ) Isso permitirá ao aluno ter condições para transitar de uma variedade da língua para outra, oral e escrita, adequando -se aos diferentes padrões e, conscientemente, fazendo uso dos recursos disponíveis para se expressar em textos próprios, bem como perceber o uso desses recursos nos textos do outro, atribuindo a eles significados nos mais variados contextos comunicativos em que estejam inseridos.
O estudo linguístico/semiótico permite, assim, formar um aluno poliglota em sua própria língua, capaz de escolher aquela mais adequada a cada momento de criação. > Letramento - Trata de “ um conjunto de práticas discursivas que envolvem os usos da escrita ” ( KLEIMAN, 2010, p. 07 ). > Alfabetização - “ Trata -se de um processo mais específico, que diz respeito à aquisição e à apropriação do sistema de escrita alfabético e ortográfico ” ( PEREIRA, 2005, p. 62 ). > Numa situação comunicativa, a interação possibilita a apropriação e a interiorização da linguagem que incidirá, progressivamente, na aprendizagem e no desenvolvimento humano ( VYGOTSKY, 2005 [ 1934 ] ).
Esses quatro eixos que estruturam o ensino e aprendizagem da língua/linguagem aportam -se em pelo menos quatro princípios fundamentais que norteiam o ensino de Língua Portuguesa : concepção sociointeracionista da linguagem, o letramento, o texto e a noção de sujeito aprendente.
A perspectiva sociointeracionista centra -se no processo de construção do conhecimento, baseado na interlocução, nos usos, portanto, na dimensão pragmática e discursiva dos estudos linguísticos, considerando não só as variedades formal e informal, oral e escrita da língua, mas a maneira como as pessoas empregam intencionalmente a língua/linguagem.
Para essa concepção, a linguagem é instrumento de ação coletiva e individual, logo de caráter social e cognitivo, que possibilita tanto o desenvolvimento intelectual, por meio da aquisição e transmissão do conhecimento, quanto à construção de uma consciência de sujeito social, que necessita aprender a conhecer o mundo e a conhecer -se como um sujeito desse mundo ; que precisa agir sobre o outro e sobre esse mundo ; que precisa compreender que atua sobre as representações de um mundo físico e social, constituído de normas e valores.
Por exemplo : a escola tem suas representações constituídas : lugar institucionalizado de ensino ( mundo físico ) ; lugar que demanda respeito, intelectualidade etc. ( mundo social ) ( BRONCKART, 2006 ).
O letramento, por sua vez, é condição de inclusão e participação do sujeito nas práticas de leitura e escrita nas diversas esferas sociais, o que implica um processo em constante desenvolvimento.
A escola, sendo, portanto, locus institucional de formação humana, tem a responsabilidade de envolver um conjunto de textos/discursos nos processos de aprendizagem do aluno, para não só alfabetizá -lo, mas torná -lo um indivíduo letrado, que não só sabe ler e escrever como também sabe usar essas habilidades para sua participação na sociedade letrada.
Assim sendo, é crucial para o êxito da educação atual que a prática pedagógica esteja fundamentada na concepção de letramento para redimensionar a compreensão do professor sobre leitura, escrita, produção textual e para conhecer as capacidades e necessidades do sujeito aprendente.
Nesse sentido, o texto é unidade básica do trabalho com a linguagem, espaço em que se concretiza todo agir de linguagem ( BRONCKART, 1999 [ 2003 ] ).
É o lugar de encontro ou desencontro de ideias, de concepções, de troca de saberes, de construção do conhecimento.
É lugar de conflito, constituído de representações ideológicas e culturais, sócio-historicamente situadas, e de valores éticos, morais e atitudinais, intrínsecos ao agir comunicativo, ou seja, ao dizer-fazer-sentir nas interações vivenciadas.
Essas formas de dizer-fazer-sentir são modeladas por organizações textuais do domínio social, que orientam as condutas humanas, não só pelas estruturas formais, mas pela situação, posição social e pela reciprocidade nas relações entre os participantes da comunicação, além do tema, estilo, são os chamados Gêneros textuais. > Gêneros textuais - “ [... ] fenômenos profundamente vinculados à vida cultural e social.
Fruto de trabalho coletivo, os gêneros contribuem para ordenar e estabilizar as atividades comunicativas do dia a dia ” ( MARCUSCHI, 2003, p. 19 ).
A concepção de sujeito aqui adotada segue a concepção sociointeracionista da linguagem, portanto é o sujeito que se constrói por meio da interação, do diálogo, do discurso ; é o sujeito que respeita a alteridade, que é responsável pelo seu dizer e compreender o outro ; um sujeito circunscrito num contexto cultural e sociopolítico e ideológico ( VOLOCHINOV, 2004 ).
É um sujeito que mantém uma relação direta com as práticas sociais, que constroem -se e reconstroem -se, conforme o contexto político, econômico, histórico e cultural.
Esse sujeito, então, tanto contribui para essas transformações como sofre as injunções por elas provocadas, o que demanda sempre revisão dos e construção de novos conhecimentos, saberes, atitudes, os quais implicarão no seu desenvolvimento cognitivo, linguístico e textual, social e afetivo.
Tudo isso concorre para a construção do sujeito social, identitário e ideologicamente firmado.
Mas qual a função da escola e especialmente a do componente Língua Portuguesa para a construção desse sujeito?
Ora, se a linguagem é o centro dinamizador da construção desse sujeito e a escola, o lugar social institucionalizado para esse fim, cabe ao ensino da língua/linguagem oportunizar experiências conceituais e vivenciais que possibilitem a interação do aluno, sujeito aprendente, com seus pares na escola, com o mundo físico e social do seu contexto imediato e/ou virtual, mediado pela representação mental ou simbólica, por meio da leitura, da produção textual, da análise linguística, empregando a oralidade, a escrita, além da multimodalidade e das multimídias. > [... ] a noção de sujeito como entidade psicossocial, sublinhando -se o caráter ativo dos sujeitos na produção mesma do social e da interação e defendendo a posição de que os sujeitos (re)produzem o social na medida em que participam ativamente da definição da situação na qual se acham engajados, e que são atores na atualização das imagens e das representações sem as quais a comunicação não poderia existir. ( KOCH, 2003, p. 15 ). > Multimodalidade - [... ] é um traço constitutivo dos gêneros.
Portanto, é no texto, materialidade dos gêneros, onde os modos ( imagem, escrita, som, música, linhas, cores, tamanho, ângulos, entonação, ritmos, efeitos visuais, melodia etc. ) são realizados.
O que faz com que um modo seja multimodal são as combinações com outros modos para criar sentidos [... ] ( DIONÍSIO, 2014, p. 42 ). > Multimídias - Sistema que combina som, imagens estáticas e animadas, vídeo, interatividade e textos, com funções educativas, entre outras. ( http://michaelis.uol.com.br/... ) Para isso, é fundamental que o ( a ) professor ( a ) organize com clareza e objetividade o seu plano de ensino, de modo a atender às necessidades do sujeito, conforme as suas etapas de maturidade biopsicointelectual, para poder impulsioná -lo, paulatinamente, e mediado pela linguagem e pelos textos, às descobertas do fazer científico, social e afetivo, acrescido dos aspectos político e ideológico, o que contribuirá para a constituição identitária desse sujeito.
Enfim, é preciso esclarecer que esses princípios estão na base do ensino de Língua Portuguesa, mas eles são, em meio a tantos outros, apenas um fio condutor para o incremento das ações didáticas que possam assegurar os Direitos de aprendizagem e desenvolvimento dos alunos no Ensino Fundamental, anos iniciais e finais.
Vale salientar que, para os 03 ( três ) primeiros anos dessa etapa ( dos quais, os dois primeiros objetivam alfabetizar os alunos, e o último, ajustar o que for necessário ao processo de alfabetização ) foram tomados por referência os direitos do Pacto Nacional pela Alfabetização na Idade Certa ( PNAIC ), que tem o propósito de garantir a todos os estudantes sua alfabetização até os oito anos de idade.
Já para o 4º ao 9º ano, os direitos de aprendizagem foram elaborados de modo a assegurar o desenvolvimento do sujeito social e identitário.
De forma geral, seja no primeiro ou no segundo nível, esses direitos foram pensados considerando o tripé início/consolidação/ampliação, em cada eixo ( oralidade, leitura, produção textual e análise linguística/semiótica ).
Por isso, os direitos de aprendizagem estão apresentados separadamente, considerando -se o ciclo da alfabetização e os anos de escolaridade do 4º ao 9º ano do Ensino Fundamental.

% matched lemmas: consubstanciar
\end{textsample}
